% Part: methods
% Chapter: induction
% Section: induction-on-N

\documentclass[../../../include/open-logic-section]{subfiles}

\begin{document}

\olfileid{mth}{ind}{inN}

\olsection{$\Nat$ પર આગમન}

સૌથી સરળ સ્વરૂપમાં આગમન બધી પ્રાકૃતિક સંખ્યાઓ માટેનાં પરિણામો
સાબિત કરવાની પદ્ધતિ છે. તેમાં એ હકીકત વપરાય છે કે $0$થી શરૂ
કરીને વારંવાર~$1$ ઉમેરતાં છેવટે દરેક પ્રાકૃતિક સંખ્યા મળે છે.
તેથી દરેક સંખ્યા માટે કંઈક સાચું છે તે સાબિત કરવા આપણે
(1)~તે $0$ માટે સાચું છે એમ સ્થાપી શકીએ અને (2)~જ્યારે તે
કોઈ સંખ્યા~$n$ માટે સાચું હોય ત્યારે પછીની સંખ્યા~$n+1$
માટે પણ સાચું છે એમ બતાવી શકીએ. "$n$માં ગુણધર્મ $P$ છે"
તેને ટૂંકમાં $P(n)$ લખીએ (અને "$k$માં ગુણધર્મ $P$ છે"ને
$P(k)$, વગેરે). તો દરેક $n \in \Nat$ માટે $P(n)$ છે એની
આગમનથી સાબિતીમાં આ બે ભાગ હોય છે:
\begin{enumerate}
\item $P(0)$ની સાબિતી, અને
\item કોઈ પણ~$k$ માટે, જો $P(k)$ તો $P(k+1)$ છે તેની સાબિતી.
\end{enumerate}
આ વાત સંપૂર્ણ સ્પષ્ટ થાય તે માટે ધારો કે (1) અને~(2) બંને
આપણી પાસે છે. ત્યારે (1) કહે છે કે $P(0)$ સાચું છે. (2) પણ
હોય તો ખાસ કરીને $P(0)$ હોય ત્યારે $P(0+1)$, એટલે $P(1)$,
થાય છે. આ "$k$ ગમે તે હોય, $P(k)$ હોય તો $P(k+1)$" એવા
સામાન્ય વિધાનમાં~$k$ માટે~$0$ મૂકવાથી મળે છે. તેથી
મોડસ પોનેન્સથી~$P(1)$ મળે છે. હવે ફરી (2)માં~$k$ માટે
$1$ મૂકતાં: $P(1)$ હોય તો~$P(2)$ મળે.\sourcecorrection{OLMD-163}{અહીં
  સ્થિર અંગ્રેજી મૂળમાં (2)ના $k$ને બદલે $n$ લખ્યું છે;
  (2)માં સર્વત્ર પરિમાણિત ચલ $k$ છે.} હમણાં જ~$P(1)$
સ્થાપ્યું હોવાથી મોડસ પોનેન્સથી~$P(2)$ મળે છે. આ રીતે
આગળ વધીએ. કોઈ પણ સંખ્યા~$n$ માટે આવું $n$~વાર કરતાં
છેવટે~$P(n)$ સુધી પહોંચીએ છીએ. આમ (1) અને~(2) મળીને
દરેક $n \in \Nat$ માટે~$P(n)$ સ્થાપે છે.

એક ઉદાહરણ જોઈએ. ધારો કે $n$ પાસા ફેંકતાં કેટલા જુદા
સરવાળા મળી શકે તે જાણવું છે. વિચિત્ર લાગે તો પણ $0$ પાસાથી
શરૂ કરીએ. એક પણ પાસો ન હોય તો ફેંકવાથી એક જ સરવાળો મળે:
એક પણ ટપકું નહીં, એટલે સરવાળો~$0$. તેથી શક્ય પરિણામોની
સંખ્યા~$1$ છે. ફક્ત એક પાસો હોય, એટલે $n=1$ હોય, તો
$1$થી~$6$ સુધીનાં છ મૂલ્યો મળે. બે પાસાથી $2$થી $12$
સુધીનો કોઈ પણ સરવાળો મળે, એટલે $11$ શક્યતાઓ. ત્રણ
પાસાથી $3$થી $18$ સુધીનાં મૂલ્યો, એટલે $16$ શક્યતાઓ
મળે. $1$, $6$, $11$, $16$: કોઈ નિયમ દેખાય છે; જવાબ
કદાચ $5n+1$ હશે? ખરેખર, $5n+1$થી વધારે મૂલ્યો શક્ય જ
નથી, કારણ કે $n$ પાસાથી મળતો લઘુતમ સરવાળો $n$ (બધા
પાસા પર $1$) અને મહત્તમ સરવાળો $6n$ (બધા પર $6$) છે,
અને તેમની વચ્ચે, છેડા સહિત, ફક્ત $5n+1$ સંખ્યાઓ છે.

\begin{thm}
  $n$ પાસાથી $n$ અને $6n$ વચ્ચેના બધા $5n+1$ શક્ય મૂલ્યો
  મેળવી શકાય છે.
\end{thm}

\begin{proof}
$P(n)$ એ "$n$ પાસા વડે $n$ અને~$6n$ વચ્ચેની કોઈ પણ
સંખ્યા મેળવી શકાય છે" એવો દાવો રાખીએ. આગમન વાપરવા
માટે આપણે સાબિત કરીએ:
\begin{enumerate}
\item \emph{આગમનનો આધાર} $P(1)$, એટલે ફક્ત એક પાસાથી
  $1$ અને $6$ વચ્ચેની દરેક સંખ્યા મેળવી શકાય છે.
\item \emph{આગમન-પગલું}: બધા $k$ માટે, જો $P(k)$ તો~$P(k+1)$.
\end{enumerate}

(1) $6$-બાજુવાળો પાસો જોઈને સાબિત થાય છે. તેને છ બાજુઓ છે
અને $1$થી~$6$ સુધીની દરેક સંખ્યા તેની કોઈ એક બાજુ પર છે.
તેથી એક જ પાસો વાપરીને $1$ અને~$6$ વચ્ચેની કોઈ પણ સંખ્યા
મેળવી શકાય છે.

(2) સાબિત કરવા આપણે શરતી વિધાનનું પૂર્વપદ, એટલે $P(k)$,
ધારીએ. આ ધારણાને \emph{આગમન-પરિકલ્પના} કહે છે. તેનો
ઉપયોગ કરીને~$P(k+1)$ સાબિત કરીએ. મુશ્કેલ ભાગ એ છે કે
$k+1$ પાસા ફેંકતાં મળતાં મૂલ્યોને $k$~પાસાથી મળતાં
મૂલ્યો અને વધારાના $k+1$મા પાસાથી મળતાં મૂલ્યોની
દૃષ્ટિએ કેવી રીતે વિચારવાં. આગમન-પરિકલ્પના વાપરવી
હોય તો આવું કરવું જ પડે.

આગમન-પરિકલ્પના મુજબ $k$~પાસાથી $k$ અને~$6k$ વચ્ચેની
કોઈ પણ સંખ્યા મેળવી શકાય છે. $(k+1)$મા પાસા પર~$1$
આવે તો કુલ સરવાળામાં $1$ ઉમેરાય. તેથી $k$~પાસા ફેંકી,
પછી $(k+1)$મા પાસા પર~$1$ મેળવીને $k+1$થી $6k+1$
સુધીનું કોઈ પણ મૂલ્ય મેળવી શકીએ. હવે શું બાકી છે?
$6k+2$થી $6k+6$ સુધીનાં મૂલ્યો. પહેલા $k$ પાસા પર
$6$ અને પછી $(k+1)$મા પાસા પર $2$થી $6$ વચ્ચેની
સંખ્યા મેળવીને આ મૂલ્યો મળે છે. આમ $k+1$ પાસાથી
$k+1$ અને $6(k+1)$ વચ્ચેની બધી સંખ્યાઓ મેળવી શકાય છે.
એટલે આગમન-પરિકલ્પના~$P(k)$ ધારીને~$P(k+1)$ સાબિત થયું.
\end{proof}

ઘણી વાર વસ્તુઓની (સંખ્યાઓ, ગણો વગેરેની) એવી શ્રેણી
વિશે કંઈક સાબિત કરવા આપણે આગમન વાપરીએ છીએ જેને
પોતાને "આગમનાત્મક રીતે" વ્યાખ્યાયિત કરી હોય, એટલે
$(n+1)$મી વસ્તુને $n$મી વસ્તુ વડે વ્યાખ્યાયિત કરી હોય.
દાખલા તરીકે,~$n$ સુધીની પ્રાકૃતિક સંખ્યાઓનો સરવાળો~$s_n$
આ રીતે વ્યાખ્યાયિત કરી શકીએ:
\begin{align*}
  s_0 & = 0\\
  s_{n+1} & = s_n + (n+1)
\end{align*}
આ વ્યાખ્યાથી મળે છે:
\begin{align*}
  s_0 & = 0,\\
  s_1 & = s_0 + 1 && = 1,\\
  s_2 & = s_1 + 2 && = 1 + 2 = 3\\
  s_3 & = s_2 + 3 && = 1 + 2 + 3 = 6, \text{ વગેરે}
\end{align*}
હવે આગમનથી સાબિત કરી શકીએ કે $s_n = n(n+1)/2$.

\begin{prop}
  $s_n = n(n+1)/2$.
\end{prop}

\begin{proof}
  આપણે (1)~$s_0 = 0\cdot(0 + 1)/2$ અને (2)~જો $s_k =
  k(k+1)/2$ તો $s_{k+1} = (k+1)(k+2)/2$ સાબિત કરવાનું છે.
  (1) સ્પષ્ટ છે. (2) માટે આગમન-પરિકલ્પના $s_k = k(k+1)/2$
  ધારીએ. તેને વાપરીને $s_{k+1} = (k+1)(k+2)/2$ બતાવવું છે.

  $s_{k+1}$ શું છે? વ્યાખ્યા મુજબ $s_{k+1} = s_k + (k+1)$.
  આગમન-પરિકલ્પના મુજબ $s_k = k(k+1)/2$. તેને અગાઉના
  સમીકરણમાં મૂકીએ તો માત્ર થોડું અપૂર્ણાંકોનું ગણિત બાકી રહે છે:
  \begin{align*}
    s_{k+1} & = \frac{k(k+1)}{2} + (k+1) = {}\\
    & = \frac{k(k+1)}{2} + \frac{2(k+1)}{2} = {}\\
    & = \frac{k(k+1) + 2(k+1)}{2} = {}\\
    & = \frac{(k+2)(k+1)}{2}.
  \end{align*}
\end{proof}

અહીં મુખ્ય પાઠ એ છે કે કોઈ આગમનાત્મક રીતે વ્યાખ્યાયિત
શ્રેણી~$a_n$ વિશે પરિણામ સાબિત કરવું હોય તો આગમન
સ્વાભાવિક પદ્ધતિ છે. અને જો શ્રેણી એવી રીતે વ્યાખ્યાયિત
ન પણ હોય (જેમ પાસા ફેંકવાના ઉદાહરણમાં), તોપણ~$k+1$નો
કિસ્સો~$k$ના કિસ્સા સાથે કોઈ રીતે જોડી શકાય તો આગમન
વાપરી શકાય છે.

\end{document}
